
Transisi menuju Energi Baru Terbarukan (EBT) pada sistem tenaga listrik menyebabkan penurunan inersia sistem secara signifikan, yang berisiko meningkatkan kerentanan terhadap ketidakstabilan frekuensi hingga pemadaman total (\textit{blackout}). Saat terjadi defisit daya, metode pertahanan darurat konvensional seperti \textit{Under Frequency Load Shedding} (UFLS) sering kali bersifat statis dan berpotensi memicu pelepasan beban yang berlebihan (\textit{overshedding}). Menanggapi permasalahan tersebut, proyek \textit{Capstone} ini mengusulkan solusi berupa sistem \textit{SCADA-Based Optimal Load Shedding and Recovery} yang mampu beroperasi secara otomatis, \textit{real-time}, dan terintegrasi. Sistem ini menggabungkan algoritma optimasi matematis dengan simulasi dinamika frekuensi dan kendali perangkat keras berbasis \textit{Programmable Logic Controller} (PLC) Schneider Electric Modicon M221. Perancangan arsitektur sistem dilakukan berdasarkan evaluasi studi literatur secara komprehensif guna memastikan jaminan optimalitas, kecepatan komputasi, serta kepatuhan terhadap standar regulasi operasi jaringan (\textit{Grid Code}) seperti Peraturan Menteri ESDM No. 3 Tahun 2007 dan IEEE Std 1547-2018.

Sistem dirancang untuk mendeteksi kondisi darurat melalui simulasi dinamika \textit{Rate of Change of Frequency} (RoCoF) dan kalkulasi defisit kapasitas daya ($\Delta P$) secara \textit{real-time}. Apabila sistem mendeteksi gangguan, keputusan pelepasan beban dieksekusi secara dinamis menggunakan algoritma \textit{Mixed-Integer Linear Programming} (MILP-DC). Algoritma ini menjamin minimasi kerugian ekonomi (\textit{Cost of Energy Not Served}) dengan tetap memprioritaskan keamanan beban kritis serta mematuhi batasan fisik operasional. Komunikasi data dua arah antara algoritma \textit{backend} Python dan PLC dijalankan secara persisten menggunakan protokol Modbus TCP/IP. Hasil komputasi kemudian diterapkan secara fisik dengan mengendalikan \textit{circuit breaker} (CB) untuk fase pemutusan darurat maupun penjadwalan pemulihan beban secara bertahap pasca-gangguan. Seluruh aktivitas sistem dicatat dalam log \textit{Sequence of Events} (SOE) dan divisualisasikan melalui \textit{Single Line Diagram} (SLD) dinamis pada antarmuka \textit{Web HMI Dashboard}, yang juga memfasilitasi operator untuk melakukan \textit{override} manual jika diperlukan. Sistem ini diharapkan mampu menjadi solusi modern dan \textit{failsafe} terhadap manajemen krisis pada jaringan sub-transmisi dan distribusi masa depan yang memiliki inersia rendah.

\textbf{Kata Kunci} --- \textit{Optimal load shedding}, SCADA, \textit{Mixed-Integer Linear Programming}, PLC Schneider, inersia rendah, \textit{hardware-in-the-loop}.